% =====================================================================
% Methods text for the mesh-based membrane topology analysis.
% Written to slot into paper/main.tex either as a new Metric 8 (after the
% voxel curvature Metric 7) or to supersede Metric 7. Style matches the
% existing "What was done / Strengths / Weaknesses / Biases / Results"
% template used throughout main.tex.
%
% To use, either % =====================================================================
% Methods text for the mesh-based membrane topology analysis.
% Written to slot into paper/main.tex either as a new Metric 8 (after the
% voxel curvature Metric 7) or to supersede Metric 7. Style matches the
% existing "What was done / Strengths / Weaknesses / Biases / Results"
% template used throughout main.tex.
%
% To use, either % =====================================================================
% Methods text for the mesh-based membrane topology analysis.
% Written to slot into paper/main.tex either as a new Metric 8 (after the
% voxel curvature Metric 7) or to supersede Metric 7. Style matches the
% existing "What was done / Strengths / Weaknesses / Biases / Results"
% template used throughout main.tex.
%
% To use, either % =====================================================================
% Methods text for the mesh-based membrane topology analysis.
% Written to slot into paper/main.tex either as a new Metric 8 (after the
% voxel curvature Metric 7) or to supersede Metric 7. Style matches the
% existing "What was done / Strengths / Weaknesses / Biases / Results"
% template used throughout main.tex.
%
% To use, either \input{methods_membrane_topology} at the appropriate
% place in main.tex, or copy the section body in.
% =====================================================================

\section*{Metric 8: Membrane Topology --- Mesh-Based Curvature, Protrusion, and Local Contact Gap}

{\small\textbf{Scripts:} \texttt{scripts/render\_membrane\_patch.py},
\texttt{scripts/regen\_galleries.py},
\texttt{ecs/geometry.py:\{signed\_mean\_curvature, smoothed\_surface\_deviation,
classify\_ecs\_facing\_vertices\}}.
\textbf{Manifest:} \texttt{figures/membranes/manifest\_3d.json}.}

\subsection*{What was done}
For each crop, a single membrane-rich cell was selected as the cell with
the greatest area of cell--ECS interface (counted as the number of cell
faces sharing a boundary with an ECS voxel; ties broken by cell index).
Crops with finer-than-16\,nm native voxel size were box-downsampled to
16\,nm isotropic prior to meshing, so that all crops feed an identical
analysis grid regardless of acquisition resolution.

The cell's binary mask was smoothed with a 3D Gaussian of physical width
$\sigma = 1.5 \cdot v_x$ nm (i.e.~$\sigma = 24$\,nm at the 16\,nm working
voxel), then surfaced by marching cubes at iso-level 0.5 (skimage
implementation). Surface vertices were placed in nm coordinates accounting
for both the cropping bounding-box offset and the smoothing pad. Sign
convention for downstream metrics is calibrated against a synthetic
convex sphere of radius 400\,nm using the same pipeline; convex membranes
return positive curvature.

Three per-vertex scalars were then computed on the resulting triangle mesh:

\begin{enumerate}[nosep, leftmargin=*]
\item \textbf{Signed mean curvature} $H$ (units: nm$^{-1}$) was computed
from the cotangent Laplacian. For each interior edge of the mesh the
cotangent weights of the two opposite angles enter a sparse Laplacian
operator $L$; the mean-curvature normal vector at vertex $i$ is
$\mathbf{H}_i \mathbf{n}_i = \frac{1}{2A_i}(L\mathbf{V})_i$, where $A_i$
is the barycentric vertex area. The magnitude $|H_i|$ is the unsigned
curvature; the sign is taken from $\operatorname{sgn}(\mathbf{H}_i \cdot \mathbf{n}_i^{\text{trimesh}})$,
calibrated so convex (membrane-into-ECS) surfaces are positive. The
radius of curvature is $R = 1/|H|$; representative values:
$H = 0.005$\,nm$^{-1}$ $\leftrightarrow$ $R = 200$\,nm (gentle bend);
$H = 0.02$\,nm$^{-1}$ $\leftrightarrow$ $R = 50$\,nm (microvillus shaft);
$H = 0.05$\,nm$^{-1}$ $\leftrightarrow$ $R = 20$\,nm (microvillus tip);
$H = 0.10$\,nm$^{-1}$ $\leftrightarrow$ $R = 10$\,nm (sharp spike).

\item \textbf{Protrusion / indentation $d$} (units: nm) was computed as
the signed normal-projected displacement of each vertex from a smoothed
reference surface generated from the same mesh. The reference was
produced by random-walk Laplacian iteration on vertex coordinates
($v_\text{new} = v - \lambda D^{-1} L v$, $\lambda = 0.5$),
with the iteration count chosen so the effective smoothing scale is
$\sigma = 60$\,nm given the mesh's mean edge length: $N \approx
\sigma^2 / (2 h^2 \lambda)$. The signed deviation is the inward-normal
projection of the original-minus-smoothed displacement vector, so $d > 0$
corresponds to a vertex protruding outward into ECS and $d < 0$ to an
inward indentation. Conceptually, $d$ answers ``does this point stick
out (or in) compared to its 60\,nm neighbourhood?'' --- distinct from
$H$ in that it is scale-aware: a 100\,nm-tall ridge gives large positive
$d$ but small $|H|$; a 5\,nm bump on flat membrane gives small $d$ but
large $|H|$.

\item \textbf{Local contact gap $g$} (units: nm) was computed as the 3D
Euclidean distance transform (EDT) of the ``not-other-cell'' indicator
field, sampled at the rounded voxel coordinate of each mesh vertex.
$g_i$ is therefore the distance from vertex $i$ to the nearest voxel
belonging to any cell other than the one being analysed.
\end{enumerate}

A membrane patch was isolated by retaining only those vertices within
one voxel (max-norm) of an ECS-labeled voxel; the per-vertex test was
dilated by two ring-neighbour iterations on the mesh edge graph to
close scattered single-vertex dropouts in the ECS-facing classifier.
Faces with at least 50\% of their vertices marked ECS-facing were
included in the patch. To avoid marching-cubes ``cap'' faces produced
where the cell touches the crop wall, faces within two voxels of any
volume face were dropped from the patch.

\paragraph{Volume-boundary handling.}
Two boundary effects are corrected in the per-channel rendering and
summary statistics. (i)~The contact gap $g_i$ is computed by an in-volume
distance transform and is an upper-bounded reading only when the
nearest cell within the crop is closer than the nearest crop wall;
formally, for every vertex $g_i^{\text{true}} \le \min(g_i^{\text{EDT}},
d_i^{\text{wall}})$ where $d_i^{\text{wall}}$ is the L$_\infty$ distance from
vertex $i$ to the nearest of the six volume faces. Vertices satisfying
$g_i^{\text{EDT}} > d_i^{\text{wall}}$ were flagged as
boundary-uncertain and dropped from the gap channel (rather than
silently clipped --- clipping just paints a low-value rim, swapping one
artifact for another). The per-patch boundary-uncertain fraction
(\texttt{gap\_bounded\_frac}, abbreviated \emph{bd-clip} in figure legends)
is reported as a quality indicator. (ii)~The cotangent Laplacian and the
60\,nm smoothing kernel both reach beyond the patch rim into the
marching-cubes cap face at the volume boundary, biasing values inward
and painting an artificial protrusion stripe along the edge of every
patch. Faces whose vertices lie within $\sigma = 60$\,nm of any volume
face were therefore dropped from the curvature and protrusion channels.
Patch geometry, per-cell statistics (face counts, ECS-facing area
fraction) are reported on the unfiltered patch so the boundary trims
affect rendering and per-channel summaries but not the geometric
denominator.

\paragraph{Per-crop summary statistics.}
For each crop's patch the following are reported in
\texttt{manifest\_3d.json}: number of patch faces, number of mesh faces
in the underlying full-cell mesh, ECS-facing area fraction, gap-channel
face count, boundary-uncertain fraction, dataset and cell ID, and the
adaptive gap colormap range (see Visualization below).

\subsection*{Strengths}
\begin{itemize}[nosep]
\item Mesh-based curvature is computed directly from the smoothed
isosurface, avoiding the finite-difference noise that limits the
voxel-based formula (Metric~7).
\item The smoothing scale ($\sigma_\text{mesh} = 1.5 v_x$ for surface
extraction, $\sigma_\text{ref} = 60$\,nm for protrusion) is set in
physical units independent of the acquisition voxel size, so HPF
(8\,nm voxel) and Chemical (4\,nm, downsampled to 16\,nm working voxel)
crops feed the same analysis at the same physical scale.
\item Three orthogonal channels (curvature, protrusion, contact gap)
capture different aspects of membrane geometry. Curvature answers
``how bent is the surface?''; protrusion answers ``how far does it
stick out compared to its 60\,nm neighbourhood?''; contact gap
answers ``how far is the nearest neighbouring cell?''. The channels
agree at sharp microvillar features and diverge over gently
undulating surfaces, providing internal cross-checks.
\item The boundary-uncertain face filter for the gap channel makes
crop-wall artifacts explicit rather than hidden: high
\texttt{bd-clip} indicates a crop whose ECS-facing patch hugs the
volume wall and whose gap channel should be interpreted cautiously.
\end{itemize}

\subsection*{Weaknesses}
\begin{itemize}[nosep]
\item Only one cell per crop is analysed (the cell with the largest
ECS-facing surface). This biases toward cells with simple convex
exposure and against cells whose ECS-facing surface is broken up by
cell-cell contacts.
\item The 16\,nm working voxel limits the resolution of the surface;
features at the 8--16\,nm scale (e.g.~very narrow microvilli in some
kidney brush borders) may be smoothed below detectability before
marching cubes is applied. Stats use 16\,nm grid uniformly across all
crops; finer-resolution Chemical crops are deliberately not given a
finer-scale advantage.
\item The 60\,nm reference scale for protrusion is a choice, not a
calibrated optimum; protrusion magnitudes shift with this scale.
\item The contact gap channel is computed only along ECS-facing
membrane and therefore does not capture cell--cell contacts where
the gap is too narrow to contain an ECS voxel. The companion
Voronoi tessellation analysis (Metric~5) reports the latter.
\item Boundary trims remove informative surface on small crops:
$\sigma = 60$\,nm corresponds to 3--4 16\,nm voxels of rim per face.
On a $1.6 \times 1.6 \times 1.6$\,\micro m kidney crop with two cells
this can drop 30--50\% of the patch from the curvature/protrusion
render even when the underlying signal is genuine.
\end{itemize}

\subsection*{Biases and artifacts}
\begin{itemize}[nosep]
\item \textbf{Resolution confound mitigation.} All crops are
downsampled to 16\,nm before analysis. Chemical crops (native 4\,nm)
therefore lose finer structure they could have resolved; HPF crops
(native 8\,nm) lose less. The choice is deliberate --- the alternative
of analysing each at native resolution makes Chemical look
artifactually sharper (higher $|H|$, larger $|d|$) for resolution
reasons unrelated to fixation.
\item \textbf{Cell selection bias.} The largest-ECS-facing-cell
criterion will favour cells that already have open ECS around them;
crops where the ``best'' cell is jammed between others may yield no
patch and be excluded. Of the 52 crops analysed, this filter produced
non-degenerate patches for all 52.
\item \textbf{Sign-of-curvature calibration.} The sign of $H$ depends
on which direction is treated as ``outward.'' We anchor the convention
to a synthetic convex sphere through the same pipeline (see
\texttt{sphere\_sign\_check}); convex surfaces return $H > 0$. Any
break of this convention (e.g.~a degenerate cell whose mesh has
inverted normals) is detectable from the sphere check and was not
observed in production crops.
\item \textbf{Disconnected gap-channel components.} On crops where a
significant fraction of vertices are boundary-uncertain, the dropped
rim can physically separate the gap-channel submesh into multiple
connected components. This is rendered honestly --- the submesh
silhouette is irregular --- but should be interpreted with the
\texttt{bd-clip} fraction.
\end{itemize}

\subsection*{Visualization}
For the public gallery (\texttt{figures/membranes/membranes\_3d.html}),
the per-vertex scalars are vertex-coloured into glTF~2.0 (\texttt{.glb})
meshes via the \texttt{matplotlib} \texttt{RdBu\_r} (curvature,
protrusion) and \texttt{viridis} (gap) colormaps and rendered with
the \texttt{model-viewer} web component. Each card displays the
patch's three channels selectable in-place and also via a global
``color all by'' selector. The curvature and protrusion colormap
limits are symmetric and set to $\pm$\,p90 of $|H|$ and $|d|$ on the
patch (the brightest reds/blues are the sharpest features
\emph{in that patch}, not a fixed absolute threshold across crops).
The gap colormap is anchored at 0--$3 \cdot 40 = 120$\,nm (so crops
dominated by close contacts share a common scale) and expanded to
p95 of the kept-vertex gap distribution when the patch genuinely
sits hundreds of nm from any in-volume neighbour, with the active
upper bound shown on each card. Each card also exposes a
Neuroglancer link wired to the underlying crop with the EM, ECS
silhouette mesh, and all cell meshes pre-loaded; this allows
visual verification of any reading against the raw FIB-SEM EM.

\subsection*{Results --- overall}
% TODO: populate from manifest_3d.json once aggregate per-tissue
% summaries are computed. Suggested form: per-tissue table of
% (median $|H|$, IQR), (median $|d|$, IQR), (median $g$, IQR), and
% the bd-clip distribution, separated by prep. The matched-region
% comparison in the following subsection is the primary analytical
% result; this paragraph is a high-level descriptive aggregate.

\subsection*{Results --- region-matched comparisons}
% TODO: populate from manifest_3d.json + crop_annotations.csv. The
% matched region groups (e.g.~Liver Bile canaliculus, Liver Space of
% Disse, Liver Hepatocyte lateral; Heart Intercalated disc, Cardiac
% interstitial; Kidney PCT brush border, PCT lateral, DCT base,
% Glomerular interstitium; Cortex Neurovascular, Synaptic,
% Mitochondrion-rich) are the unit of comparison. Suggested form for
% each region group: paired per-channel (Chem vs HPF) per-vertex
% distributions, summarized as median + IQR, with a small inset
% mesh showing the most-representative crop.
 at the appropriate
% place in main.tex, or copy the section body in.
% =====================================================================

\section*{Metric 8: Membrane Topology --- Mesh-Based Curvature, Protrusion, and Local Contact Gap}

{\small\textbf{Scripts:} \texttt{scripts/render\_membrane\_patch.py},
\texttt{scripts/regen\_galleries.py},
\texttt{ecs/geometry.py:\{signed\_mean\_curvature, smoothed\_surface\_deviation,
classify\_ecs\_facing\_vertices\}}.
\textbf{Manifest:} \texttt{figures/membranes/manifest\_3d.json}.}

\subsection*{What was done}
For each crop, a single membrane-rich cell was selected as the cell with
the greatest area of cell--ECS interface (counted as the number of cell
faces sharing a boundary with an ECS voxel; ties broken by cell index).
Crops with finer-than-16\,nm native voxel size were box-downsampled to
16\,nm isotropic prior to meshing, so that all crops feed an identical
analysis grid regardless of acquisition resolution.

The cell's binary mask was smoothed with a 3D Gaussian of physical width
$\sigma = 1.5 \cdot v_x$ nm (i.e.~$\sigma = 24$\,nm at the 16\,nm working
voxel), then surfaced by marching cubes at iso-level 0.5 (skimage
implementation). Surface vertices were placed in nm coordinates accounting
for both the cropping bounding-box offset and the smoothing pad. Sign
convention for downstream metrics is calibrated against a synthetic
convex sphere of radius 400\,nm using the same pipeline; convex membranes
return positive curvature.

Three per-vertex scalars were then computed on the resulting triangle mesh:

\begin{enumerate}[nosep, leftmargin=*]
\item \textbf{Signed mean curvature} $H$ (units: nm$^{-1}$) was computed
from the cotangent Laplacian. For each interior edge of the mesh the
cotangent weights of the two opposite angles enter a sparse Laplacian
operator $L$; the mean-curvature normal vector at vertex $i$ is
$\mathbf{H}_i \mathbf{n}_i = \frac{1}{2A_i}(L\mathbf{V})_i$, where $A_i$
is the barycentric vertex area. The magnitude $|H_i|$ is the unsigned
curvature; the sign is taken from $\operatorname{sgn}(\mathbf{H}_i \cdot \mathbf{n}_i^{\text{trimesh}})$,
calibrated so convex (membrane-into-ECS) surfaces are positive. The
radius of curvature is $R = 1/|H|$; representative values:
$H = 0.005$\,nm$^{-1}$ $\leftrightarrow$ $R = 200$\,nm (gentle bend);
$H = 0.02$\,nm$^{-1}$ $\leftrightarrow$ $R = 50$\,nm (microvillus shaft);
$H = 0.05$\,nm$^{-1}$ $\leftrightarrow$ $R = 20$\,nm (microvillus tip);
$H = 0.10$\,nm$^{-1}$ $\leftrightarrow$ $R = 10$\,nm (sharp spike).

\item \textbf{Protrusion / indentation $d$} (units: nm) was computed as
the signed normal-projected displacement of each vertex from a smoothed
reference surface generated from the same mesh. The reference was
produced by random-walk Laplacian iteration on vertex coordinates
($v_\text{new} = v - \lambda D^{-1} L v$, $\lambda = 0.5$),
with the iteration count chosen so the effective smoothing scale is
$\sigma = 60$\,nm given the mesh's mean edge length: $N \approx
\sigma^2 / (2 h^2 \lambda)$. The signed deviation is the inward-normal
projection of the original-minus-smoothed displacement vector, so $d > 0$
corresponds to a vertex protruding outward into ECS and $d < 0$ to an
inward indentation. Conceptually, $d$ answers ``does this point stick
out (or in) compared to its 60\,nm neighbourhood?'' --- distinct from
$H$ in that it is scale-aware: a 100\,nm-tall ridge gives large positive
$d$ but small $|H|$; a 5\,nm bump on flat membrane gives small $d$ but
large $|H|$.

\item \textbf{Local contact gap $g$} (units: nm) was computed as the 3D
Euclidean distance transform (EDT) of the ``not-other-cell'' indicator
field, sampled at the rounded voxel coordinate of each mesh vertex.
$g_i$ is therefore the distance from vertex $i$ to the nearest voxel
belonging to any cell other than the one being analysed.
\end{enumerate}

A membrane patch was isolated by retaining only those vertices within
one voxel (max-norm) of an ECS-labeled voxel; the per-vertex test was
dilated by two ring-neighbour iterations on the mesh edge graph to
close scattered single-vertex dropouts in the ECS-facing classifier.
Faces with at least 50\% of their vertices marked ECS-facing were
included in the patch. To avoid marching-cubes ``cap'' faces produced
where the cell touches the crop wall, faces within two voxels of any
volume face were dropped from the patch.

\paragraph{Volume-boundary handling.}
Two boundary effects are corrected in the per-channel rendering and
summary statistics. (i)~The contact gap $g_i$ is computed by an in-volume
distance transform and is an upper-bounded reading only when the
nearest cell within the crop is closer than the nearest crop wall;
formally, for every vertex $g_i^{\text{true}} \le \min(g_i^{\text{EDT}},
d_i^{\text{wall}})$ where $d_i^{\text{wall}}$ is the L$_\infty$ distance from
vertex $i$ to the nearest of the six volume faces. Vertices satisfying
$g_i^{\text{EDT}} > d_i^{\text{wall}}$ were flagged as
boundary-uncertain and dropped from the gap channel (rather than
silently clipped --- clipping just paints a low-value rim, swapping one
artifact for another). The per-patch boundary-uncertain fraction
(\texttt{gap\_bounded\_frac}, abbreviated \emph{bd-clip} in figure legends)
is reported as a quality indicator. (ii)~The cotangent Laplacian and the
60\,nm smoothing kernel both reach beyond the patch rim into the
marching-cubes cap face at the volume boundary, biasing values inward
and painting an artificial protrusion stripe along the edge of every
patch. Faces whose vertices lie within $\sigma = 60$\,nm of any volume
face were therefore dropped from the curvature and protrusion channels.
Patch geometry, per-cell statistics (face counts, ECS-facing area
fraction) are reported on the unfiltered patch so the boundary trims
affect rendering and per-channel summaries but not the geometric
denominator.

\paragraph{Per-crop summary statistics.}
For each crop's patch the following are reported in
\texttt{manifest\_3d.json}: number of patch faces, number of mesh faces
in the underlying full-cell mesh, ECS-facing area fraction, gap-channel
face count, boundary-uncertain fraction, dataset and cell ID, and the
adaptive gap colormap range (see Visualization below).

\subsection*{Strengths}
\begin{itemize}[nosep]
\item Mesh-based curvature is computed directly from the smoothed
isosurface, avoiding the finite-difference noise that limits the
voxel-based formula (Metric~7).
\item The smoothing scale ($\sigma_\text{mesh} = 1.5 v_x$ for surface
extraction, $\sigma_\text{ref} = 60$\,nm for protrusion) is set in
physical units independent of the acquisition voxel size, so HPF
(8\,nm voxel) and Chemical (4\,nm, downsampled to 16\,nm working voxel)
crops feed the same analysis at the same physical scale.
\item Three orthogonal channels (curvature, protrusion, contact gap)
capture different aspects of membrane geometry. Curvature answers
``how bent is the surface?''; protrusion answers ``how far does it
stick out compared to its 60\,nm neighbourhood?''; contact gap
answers ``how far is the nearest neighbouring cell?''. The channels
agree at sharp microvillar features and diverge over gently
undulating surfaces, providing internal cross-checks.
\item The boundary-uncertain face filter for the gap channel makes
crop-wall artifacts explicit rather than hidden: high
\texttt{bd-clip} indicates a crop whose ECS-facing patch hugs the
volume wall and whose gap channel should be interpreted cautiously.
\end{itemize}

\subsection*{Weaknesses}
\begin{itemize}[nosep]
\item Only one cell per crop is analysed (the cell with the largest
ECS-facing surface). This biases toward cells with simple convex
exposure and against cells whose ECS-facing surface is broken up by
cell-cell contacts.
\item The 16\,nm working voxel limits the resolution of the surface;
features at the 8--16\,nm scale (e.g.~very narrow microvilli in some
kidney brush borders) may be smoothed below detectability before
marching cubes is applied. Stats use 16\,nm grid uniformly across all
crops; finer-resolution Chemical crops are deliberately not given a
finer-scale advantage.
\item The 60\,nm reference scale for protrusion is a choice, not a
calibrated optimum; protrusion magnitudes shift with this scale.
\item The contact gap channel is computed only along ECS-facing
membrane and therefore does not capture cell--cell contacts where
the gap is too narrow to contain an ECS voxel. The companion
Voronoi tessellation analysis (Metric~5) reports the latter.
\item Boundary trims remove informative surface on small crops:
$\sigma = 60$\,nm corresponds to 3--4 16\,nm voxels of rim per face.
On a $1.6 \times 1.6 \times 1.6$\,\micro m kidney crop with two cells
this can drop 30--50\% of the patch from the curvature/protrusion
render even when the underlying signal is genuine.
\end{itemize}

\subsection*{Biases and artifacts}
\begin{itemize}[nosep]
\item \textbf{Resolution confound mitigation.} All crops are
downsampled to 16\,nm before analysis. Chemical crops (native 4\,nm)
therefore lose finer structure they could have resolved; HPF crops
(native 8\,nm) lose less. The choice is deliberate --- the alternative
of analysing each at native resolution makes Chemical look
artifactually sharper (higher $|H|$, larger $|d|$) for resolution
reasons unrelated to fixation.
\item \textbf{Cell selection bias.} The largest-ECS-facing-cell
criterion will favour cells that already have open ECS around them;
crops where the ``best'' cell is jammed between others may yield no
patch and be excluded. Of the 52 crops analysed, this filter produced
non-degenerate patches for all 52.
\item \textbf{Sign-of-curvature calibration.} The sign of $H$ depends
on which direction is treated as ``outward.'' We anchor the convention
to a synthetic convex sphere through the same pipeline (see
\texttt{sphere\_sign\_check}); convex surfaces return $H > 0$. Any
break of this convention (e.g.~a degenerate cell whose mesh has
inverted normals) is detectable from the sphere check and was not
observed in production crops.
\item \textbf{Disconnected gap-channel components.} On crops where a
significant fraction of vertices are boundary-uncertain, the dropped
rim can physically separate the gap-channel submesh into multiple
connected components. This is rendered honestly --- the submesh
silhouette is irregular --- but should be interpreted with the
\texttt{bd-clip} fraction.
\end{itemize}

\subsection*{Visualization}
For the public gallery (\texttt{figures/membranes/membranes\_3d.html}),
the per-vertex scalars are vertex-coloured into glTF~2.0 (\texttt{.glb})
meshes via the \texttt{matplotlib} \texttt{RdBu\_r} (curvature,
protrusion) and \texttt{viridis} (gap) colormaps and rendered with
the \texttt{model-viewer} web component. Each card displays the
patch's three channels selectable in-place and also via a global
``color all by'' selector. The curvature and protrusion colormap
limits are symmetric and set to $\pm$\,p90 of $|H|$ and $|d|$ on the
patch (the brightest reds/blues are the sharpest features
\emph{in that patch}, not a fixed absolute threshold across crops).
The gap colormap is anchored at 0--$3 \cdot 40 = 120$\,nm (so crops
dominated by close contacts share a common scale) and expanded to
p95 of the kept-vertex gap distribution when the patch genuinely
sits hundreds of nm from any in-volume neighbour, with the active
upper bound shown on each card. Each card also exposes a
Neuroglancer link wired to the underlying crop with the EM, ECS
silhouette mesh, and all cell meshes pre-loaded; this allows
visual verification of any reading against the raw FIB-SEM EM.

\subsection*{Results --- overall}
% TODO: populate from manifest_3d.json once aggregate per-tissue
% summaries are computed. Suggested form: per-tissue table of
% (median $|H|$, IQR), (median $|d|$, IQR), (median $g$, IQR), and
% the bd-clip distribution, separated by prep. The matched-region
% comparison in the following subsection is the primary analytical
% result; this paragraph is a high-level descriptive aggregate.

\subsection*{Results --- region-matched comparisons}
% TODO: populate from manifest_3d.json + crop_annotations.csv. The
% matched region groups (e.g.~Liver Bile canaliculus, Liver Space of
% Disse, Liver Hepatocyte lateral; Heart Intercalated disc, Cardiac
% interstitial; Kidney PCT brush border, PCT lateral, DCT base,
% Glomerular interstitium; Cortex Neurovascular, Synaptic,
% Mitochondrion-rich) are the unit of comparison. Suggested form for
% each region group: paired per-channel (Chem vs HPF) per-vertex
% distributions, summarized as median + IQR, with a small inset
% mesh showing the most-representative crop.
 at the appropriate
% place in main.tex, or copy the section body in.
% =====================================================================

\section*{Metric 8: Membrane Topology --- Mesh-Based Curvature, Protrusion, and Local Contact Gap}

{\small\textbf{Scripts:} \texttt{scripts/render\_membrane\_patch.py},
\texttt{scripts/regen\_galleries.py},
\texttt{ecs/geometry.py:\{signed\_mean\_curvature, smoothed\_surface\_deviation,
classify\_ecs\_facing\_vertices\}}.
\textbf{Manifest:} \texttt{figures/membranes/manifest\_3d.json}.}

\subsection*{What was done}
For each crop, a single membrane-rich cell was selected as the cell with
the greatest area of cell--ECS interface (counted as the number of cell
faces sharing a boundary with an ECS voxel; ties broken by cell index).
Crops with finer-than-16\,nm native voxel size were box-downsampled to
16\,nm isotropic prior to meshing, so that all crops feed an identical
analysis grid regardless of acquisition resolution.

The cell's binary mask was smoothed with a 3D Gaussian of physical width
$\sigma = 1.5 \cdot v_x$ nm (i.e.~$\sigma = 24$\,nm at the 16\,nm working
voxel), then surfaced by marching cubes at iso-level 0.5 (skimage
implementation). Surface vertices were placed in nm coordinates accounting
for both the cropping bounding-box offset and the smoothing pad. Sign
convention for downstream metrics is calibrated against a synthetic
convex sphere of radius 400\,nm using the same pipeline; convex membranes
return positive curvature.

Three per-vertex scalars were then computed on the resulting triangle mesh:

\begin{enumerate}[nosep, leftmargin=*]
\item \textbf{Signed mean curvature} $H$ (units: nm$^{-1}$) was computed
from the cotangent Laplacian. For each interior edge of the mesh the
cotangent weights of the two opposite angles enter a sparse Laplacian
operator $L$; the mean-curvature normal vector at vertex $i$ is
$\mathbf{H}_i \mathbf{n}_i = \frac{1}{2A_i}(L\mathbf{V})_i$, where $A_i$
is the barycentric vertex area. The magnitude $|H_i|$ is the unsigned
curvature; the sign is taken from $\operatorname{sgn}(\mathbf{H}_i \cdot \mathbf{n}_i^{\text{trimesh}})$,
calibrated so convex (membrane-into-ECS) surfaces are positive. The
radius of curvature is $R = 1/|H|$; representative values:
$H = 0.005$\,nm$^{-1}$ $\leftrightarrow$ $R = 200$\,nm (gentle bend);
$H = 0.02$\,nm$^{-1}$ $\leftrightarrow$ $R = 50$\,nm (microvillus shaft);
$H = 0.05$\,nm$^{-1}$ $\leftrightarrow$ $R = 20$\,nm (microvillus tip);
$H = 0.10$\,nm$^{-1}$ $\leftrightarrow$ $R = 10$\,nm (sharp spike).

\item \textbf{Protrusion / indentation $d$} (units: nm) was computed as
the signed normal-projected displacement of each vertex from a smoothed
reference surface generated from the same mesh. The reference was
produced by random-walk Laplacian iteration on vertex coordinates
($v_\text{new} = v - \lambda D^{-1} L v$, $\lambda = 0.5$),
with the iteration count chosen so the effective smoothing scale is
$\sigma = 60$\,nm given the mesh's mean edge length: $N \approx
\sigma^2 / (2 h^2 \lambda)$. The signed deviation is the inward-normal
projection of the original-minus-smoothed displacement vector, so $d > 0$
corresponds to a vertex protruding outward into ECS and $d < 0$ to an
inward indentation. Conceptually, $d$ answers ``does this point stick
out (or in) compared to its 60\,nm neighbourhood?'' --- distinct from
$H$ in that it is scale-aware: a 100\,nm-tall ridge gives large positive
$d$ but small $|H|$; a 5\,nm bump on flat membrane gives small $d$ but
large $|H|$.

\item \textbf{Local contact gap $g$} (units: nm) was computed as the 3D
Euclidean distance transform (EDT) of the ``not-other-cell'' indicator
field, sampled at the rounded voxel coordinate of each mesh vertex.
$g_i$ is therefore the distance from vertex $i$ to the nearest voxel
belonging to any cell other than the one being analysed.
\end{enumerate}

A membrane patch was isolated by retaining only those vertices within
one voxel (max-norm) of an ECS-labeled voxel; the per-vertex test was
dilated by two ring-neighbour iterations on the mesh edge graph to
close scattered single-vertex dropouts in the ECS-facing classifier.
Faces with at least 50\% of their vertices marked ECS-facing were
included in the patch. To avoid marching-cubes ``cap'' faces produced
where the cell touches the crop wall, faces within two voxels of any
volume face were dropped from the patch.

\paragraph{Volume-boundary handling.}
Two boundary effects are corrected in the per-channel rendering and
summary statistics. (i)~The contact gap $g_i$ is computed by an in-volume
distance transform and is an upper-bounded reading only when the
nearest cell within the crop is closer than the nearest crop wall;
formally, for every vertex $g_i^{\text{true}} \le \min(g_i^{\text{EDT}},
d_i^{\text{wall}})$ where $d_i^{\text{wall}}$ is the L$_\infty$ distance from
vertex $i$ to the nearest of the six volume faces. Vertices satisfying
$g_i^{\text{EDT}} > d_i^{\text{wall}}$ were flagged as
boundary-uncertain and dropped from the gap channel (rather than
silently clipped --- clipping just paints a low-value rim, swapping one
artifact for another). The per-patch boundary-uncertain fraction
(\texttt{gap\_bounded\_frac}, abbreviated \emph{bd-clip} in figure legends)
is reported as a quality indicator. (ii)~The cotangent Laplacian and the
60\,nm smoothing kernel both reach beyond the patch rim into the
marching-cubes cap face at the volume boundary, biasing values inward
and painting an artificial protrusion stripe along the edge of every
patch. Faces whose vertices lie within $\sigma = 60$\,nm of any volume
face were therefore dropped from the curvature and protrusion channels.
Patch geometry, per-cell statistics (face counts, ECS-facing area
fraction) are reported on the unfiltered patch so the boundary trims
affect rendering and per-channel summaries but not the geometric
denominator.

\paragraph{Per-crop summary statistics.}
For each crop's patch the following are reported in
\texttt{manifest\_3d.json}: number of patch faces, number of mesh faces
in the underlying full-cell mesh, ECS-facing area fraction, gap-channel
face count, boundary-uncertain fraction, dataset and cell ID, and the
adaptive gap colormap range (see Visualization below).

\subsection*{Strengths}
\begin{itemize}[nosep]
\item Mesh-based curvature is computed directly from the smoothed
isosurface, avoiding the finite-difference noise that limits the
voxel-based formula (Metric~7).
\item The smoothing scale ($\sigma_\text{mesh} = 1.5 v_x$ for surface
extraction, $\sigma_\text{ref} = 60$\,nm for protrusion) is set in
physical units independent of the acquisition voxel size, so HPF
(8\,nm voxel) and Chemical (4\,nm, downsampled to 16\,nm working voxel)
crops feed the same analysis at the same physical scale.
\item Three orthogonal channels (curvature, protrusion, contact gap)
capture different aspects of membrane geometry. Curvature answers
``how bent is the surface?''; protrusion answers ``how far does it
stick out compared to its 60\,nm neighbourhood?''; contact gap
answers ``how far is the nearest neighbouring cell?''. The channels
agree at sharp microvillar features and diverge over gently
undulating surfaces, providing internal cross-checks.
\item The boundary-uncertain face filter for the gap channel makes
crop-wall artifacts explicit rather than hidden: high
\texttt{bd-clip} indicates a crop whose ECS-facing patch hugs the
volume wall and whose gap channel should be interpreted cautiously.
\end{itemize}

\subsection*{Weaknesses}
\begin{itemize}[nosep]
\item Only one cell per crop is analysed (the cell with the largest
ECS-facing surface). This biases toward cells with simple convex
exposure and against cells whose ECS-facing surface is broken up by
cell-cell contacts.
\item The 16\,nm working voxel limits the resolution of the surface;
features at the 8--16\,nm scale (e.g.~very narrow microvilli in some
kidney brush borders) may be smoothed below detectability before
marching cubes is applied. Stats use 16\,nm grid uniformly across all
crops; finer-resolution Chemical crops are deliberately not given a
finer-scale advantage.
\item The 60\,nm reference scale for protrusion is a choice, not a
calibrated optimum; protrusion magnitudes shift with this scale.
\item The contact gap channel is computed only along ECS-facing
membrane and therefore does not capture cell--cell contacts where
the gap is too narrow to contain an ECS voxel. The companion
Voronoi tessellation analysis (Metric~5) reports the latter.
\item Boundary trims remove informative surface on small crops:
$\sigma = 60$\,nm corresponds to 3--4 16\,nm voxels of rim per face.
On a $1.6 \times 1.6 \times 1.6$\,\micro m kidney crop with two cells
this can drop 30--50\% of the patch from the curvature/protrusion
render even when the underlying signal is genuine.
\end{itemize}

\subsection*{Biases and artifacts}
\begin{itemize}[nosep]
\item \textbf{Resolution confound mitigation.} All crops are
downsampled to 16\,nm before analysis. Chemical crops (native 4\,nm)
therefore lose finer structure they could have resolved; HPF crops
(native 8\,nm) lose less. The choice is deliberate --- the alternative
of analysing each at native resolution makes Chemical look
artifactually sharper (higher $|H|$, larger $|d|$) for resolution
reasons unrelated to fixation.
\item \textbf{Cell selection bias.} The largest-ECS-facing-cell
criterion will favour cells that already have open ECS around them;
crops where the ``best'' cell is jammed between others may yield no
patch and be excluded. Of the 52 crops analysed, this filter produced
non-degenerate patches for all 52.
\item \textbf{Sign-of-curvature calibration.} The sign of $H$ depends
on which direction is treated as ``outward.'' We anchor the convention
to a synthetic convex sphere through the same pipeline (see
\texttt{sphere\_sign\_check}); convex surfaces return $H > 0$. Any
break of this convention (e.g.~a degenerate cell whose mesh has
inverted normals) is detectable from the sphere check and was not
observed in production crops.
\item \textbf{Disconnected gap-channel components.} On crops where a
significant fraction of vertices are boundary-uncertain, the dropped
rim can physically separate the gap-channel submesh into multiple
connected components. This is rendered honestly --- the submesh
silhouette is irregular --- but should be interpreted with the
\texttt{bd-clip} fraction.
\end{itemize}

\subsection*{Visualization}
For the public gallery (\texttt{figures/membranes/membranes\_3d.html}),
the per-vertex scalars are vertex-coloured into glTF~2.0 (\texttt{.glb})
meshes via the \texttt{matplotlib} \texttt{RdBu\_r} (curvature,
protrusion) and \texttt{viridis} (gap) colormaps and rendered with
the \texttt{model-viewer} web component. Each card displays the
patch's three channels selectable in-place and also via a global
``color all by'' selector. The curvature and protrusion colormap
limits are symmetric and set to $\pm$\,p90 of $|H|$ and $|d|$ on the
patch (the brightest reds/blues are the sharpest features
\emph{in that patch}, not a fixed absolute threshold across crops).
The gap colormap is anchored at 0--$3 \cdot 40 = 120$\,nm (so crops
dominated by close contacts share a common scale) and expanded to
p95 of the kept-vertex gap distribution when the patch genuinely
sits hundreds of nm from any in-volume neighbour, with the active
upper bound shown on each card. Each card also exposes a
Neuroglancer link wired to the underlying crop with the EM, ECS
silhouette mesh, and all cell meshes pre-loaded; this allows
visual verification of any reading against the raw FIB-SEM EM.

\subsection*{Results --- overall}
% TODO: populate from manifest_3d.json once aggregate per-tissue
% summaries are computed. Suggested form: per-tissue table of
% (median $|H|$, IQR), (median $|d|$, IQR), (median $g$, IQR), and
% the bd-clip distribution, separated by prep. The matched-region
% comparison in the following subsection is the primary analytical
% result; this paragraph is a high-level descriptive aggregate.

\subsection*{Results --- region-matched comparisons}
% TODO: populate from manifest_3d.json + crop_annotations.csv. The
% matched region groups (e.g.~Liver Bile canaliculus, Liver Space of
% Disse, Liver Hepatocyte lateral; Heart Intercalated disc, Cardiac
% interstitial; Kidney PCT brush border, PCT lateral, DCT base,
% Glomerular interstitium; Cortex Neurovascular, Synaptic,
% Mitochondrion-rich) are the unit of comparison. Suggested form for
% each region group: paired per-channel (Chem vs HPF) per-vertex
% distributions, summarized as median + IQR, with a small inset
% mesh showing the most-representative crop.
 at the appropriate
% place in main.tex, or copy the section body in.
% =====================================================================

\section*{Metric 8: Membrane Topology --- Mesh-Based Curvature, Protrusion, and Local Contact Gap}

{\small\textbf{Scripts:} \texttt{scripts/render\_membrane\_patch.py},
\texttt{scripts/regen\_galleries.py},
\texttt{ecs/geometry.py:\{signed\_mean\_curvature, smoothed\_surface\_deviation,
classify\_ecs\_facing\_vertices\}}.
\textbf{Manifest:} \texttt{figures/membranes/manifest\_3d.json}.}

\subsection*{What was done}
For each crop, a single membrane-rich cell was selected as the cell with
the greatest area of cell--ECS interface (counted as the number of cell
faces sharing a boundary with an ECS voxel; ties broken by cell index).
Crops with finer-than-16\,nm native voxel size were box-downsampled to
16\,nm isotropic prior to meshing, so that all crops feed an identical
analysis grid regardless of acquisition resolution.

The cell's binary mask was smoothed with a 3D Gaussian of physical width
$\sigma = 1.5 \cdot v_x$ nm (i.e.~$\sigma = 24$\,nm at the 16\,nm working
voxel), then surfaced by marching cubes at iso-level 0.5 (skimage
implementation). Surface vertices were placed in nm coordinates accounting
for both the cropping bounding-box offset and the smoothing pad. Sign
convention for downstream metrics is calibrated against a synthetic
convex sphere of radius 400\,nm using the same pipeline; convex membranes
return positive curvature.

Three per-vertex scalars were then computed on the resulting triangle mesh:

\begin{enumerate}[nosep, leftmargin=*]
\item \textbf{Signed mean curvature} $H$ (units: nm$^{-1}$) was computed
from the cotangent Laplacian. For each interior edge of the mesh the
cotangent weights of the two opposite angles enter a sparse Laplacian
operator $L$; the mean-curvature normal vector at vertex $i$ is
$\mathbf{H}_i \mathbf{n}_i = \frac{1}{2A_i}(L\mathbf{V})_i$, where $A_i$
is the barycentric vertex area. The magnitude $|H_i|$ is the unsigned
curvature; the sign is taken from $\operatorname{sgn}(\mathbf{H}_i \cdot \mathbf{n}_i^{\text{trimesh}})$,
calibrated so convex (membrane-into-ECS) surfaces are positive. The
radius of curvature is $R = 1/|H|$; representative values:
$H = 0.005$\,nm$^{-1}$ $\leftrightarrow$ $R = 200$\,nm (gentle bend);
$H = 0.02$\,nm$^{-1}$ $\leftrightarrow$ $R = 50$\,nm (microvillus shaft);
$H = 0.05$\,nm$^{-1}$ $\leftrightarrow$ $R = 20$\,nm (microvillus tip);
$H = 0.10$\,nm$^{-1}$ $\leftrightarrow$ $R = 10$\,nm (sharp spike).

\item \textbf{Protrusion / indentation $d$} (units: nm) was computed as
the signed normal-projected displacement of each vertex from a smoothed
reference surface generated from the same mesh. The reference was
produced by random-walk Laplacian iteration on vertex coordinates
($v_\text{new} = v - \lambda D^{-1} L v$, $\lambda = 0.5$),
with the iteration count chosen so the effective smoothing scale is
$\sigma = 60$\,nm given the mesh's mean edge length: $N \approx
\sigma^2 / (2 h^2 \lambda)$. The signed deviation is the inward-normal
projection of the original-minus-smoothed displacement vector, so $d > 0$
corresponds to a vertex protruding outward into ECS and $d < 0$ to an
inward indentation. Conceptually, $d$ answers ``does this point stick
out (or in) compared to its 60\,nm neighbourhood?'' --- distinct from
$H$ in that it is scale-aware: a 100\,nm-tall ridge gives large positive
$d$ but small $|H|$; a 5\,nm bump on flat membrane gives small $d$ but
large $|H|$.

\item \textbf{Local contact gap $g$} (units: nm) was computed as the 3D
Euclidean distance transform (EDT) of the ``not-other-cell'' indicator
field, sampled at the rounded voxel coordinate of each mesh vertex.
$g_i$ is therefore the distance from vertex $i$ to the nearest voxel
belonging to any cell other than the one being analysed.
\end{enumerate}

A membrane patch was isolated by retaining only those vertices within
one voxel (max-norm) of an ECS-labeled voxel; the per-vertex test was
dilated by two ring-neighbour iterations on the mesh edge graph to
close scattered single-vertex dropouts in the ECS-facing classifier.
Faces with at least 50\% of their vertices marked ECS-facing were
included in the patch. To avoid marching-cubes ``cap'' faces produced
where the cell touches the crop wall, faces within two voxels of any
volume face were dropped from the patch.

\paragraph{Volume-boundary handling.}
Two boundary effects are corrected in the per-channel rendering and
summary statistics. (i)~The contact gap $g_i$ is computed by an in-volume
distance transform and is an upper-bounded reading only when the
nearest cell within the crop is closer than the nearest crop wall;
formally, for every vertex $g_i^{\text{true}} \le \min(g_i^{\text{EDT}},
d_i^{\text{wall}})$ where $d_i^{\text{wall}}$ is the L$_\infty$ distance from
vertex $i$ to the nearest of the six volume faces. Vertices satisfying
$g_i^{\text{EDT}} > d_i^{\text{wall}}$ were flagged as
boundary-uncertain and dropped from the gap channel (rather than
silently clipped --- clipping just paints a low-value rim, swapping one
artifact for another). The per-patch boundary-uncertain fraction
(\texttt{gap\_bounded\_frac}, abbreviated \emph{bd-clip} in figure legends)
is reported as a quality indicator. (ii)~The cotangent Laplacian and the
60\,nm smoothing kernel both reach beyond the patch rim into the
marching-cubes cap face at the volume boundary, biasing values inward
and painting an artificial protrusion stripe along the edge of every
patch. Faces whose vertices lie within $\sigma = 60$\,nm of any volume
face were therefore dropped from the curvature and protrusion channels.
Patch geometry, per-cell statistics (face counts, ECS-facing area
fraction) are reported on the unfiltered patch so the boundary trims
affect rendering and per-channel summaries but not the geometric
denominator.

\paragraph{Per-crop summary statistics.}
For each crop's patch the following are reported in
\texttt{manifest\_3d.json}: number of patch faces, number of mesh faces
in the underlying full-cell mesh, ECS-facing area fraction, gap-channel
face count, boundary-uncertain fraction, dataset and cell ID, and the
adaptive gap colormap range (see Visualization below).

\subsection*{Strengths}
\begin{itemize}[nosep]
\item Mesh-based curvature is computed directly from the smoothed
isosurface, avoiding the finite-difference noise that limits the
voxel-based formula (Metric~7).
\item The smoothing scale ($\sigma_\text{mesh} = 1.5 v_x$ for surface
extraction, $\sigma_\text{ref} = 60$\,nm for protrusion) is set in
physical units independent of the acquisition voxel size, so HPF
(8\,nm voxel) and Chemical (4\,nm, downsampled to 16\,nm working voxel)
crops feed the same analysis at the same physical scale.
\item Three orthogonal channels (curvature, protrusion, contact gap)
capture different aspects of membrane geometry. Curvature answers
``how bent is the surface?''; protrusion answers ``how far does it
stick out compared to its 60\,nm neighbourhood?''; contact gap
answers ``how far is the nearest neighbouring cell?''. The channels
agree at sharp microvillar features and diverge over gently
undulating surfaces, providing internal cross-checks.
\item The boundary-uncertain face filter for the gap channel makes
crop-wall artifacts explicit rather than hidden: high
\texttt{bd-clip} indicates a crop whose ECS-facing patch hugs the
volume wall and whose gap channel should be interpreted cautiously.
\end{itemize}

\subsection*{Weaknesses}
\begin{itemize}[nosep]
\item Only one cell per crop is analysed (the cell with the largest
ECS-facing surface). This biases toward cells with simple convex
exposure and against cells whose ECS-facing surface is broken up by
cell-cell contacts.
\item The 16\,nm working voxel limits the resolution of the surface;
features at the 8--16\,nm scale (e.g.~very narrow microvilli in some
kidney brush borders) may be smoothed below detectability before
marching cubes is applied. Stats use 16\,nm grid uniformly across all
crops; finer-resolution Chemical crops are deliberately not given a
finer-scale advantage.
\item The 60\,nm reference scale for protrusion is a choice, not a
calibrated optimum; protrusion magnitudes shift with this scale.
\item The contact gap channel is computed only along ECS-facing
membrane and therefore does not capture cell--cell contacts where
the gap is too narrow to contain an ECS voxel. The companion
Voronoi tessellation analysis (Metric~5) reports the latter.
\item Boundary trims remove informative surface on small crops:
$\sigma = 60$\,nm corresponds to 3--4 16\,nm voxels of rim per face.
On a $1.6 \times 1.6 \times 1.6$\,\micro m kidney crop with two cells
this can drop 30--50\% of the patch from the curvature/protrusion
render even when the underlying signal is genuine.
\end{itemize}

\subsection*{Biases and artifacts}
\begin{itemize}[nosep]
\item \textbf{Resolution confound mitigation.} All crops are
downsampled to 16\,nm before analysis. Chemical crops (native 4\,nm)
therefore lose finer structure they could have resolved; HPF crops
(native 8\,nm) lose less. The choice is deliberate --- the alternative
of analysing each at native resolution makes Chemical look
artifactually sharper (higher $|H|$, larger $|d|$) for resolution
reasons unrelated to fixation.
\item \textbf{Cell selection bias.} The largest-ECS-facing-cell
criterion will favour cells that already have open ECS around them;
crops where the ``best'' cell is jammed between others may yield no
patch and be excluded. Of the 52 crops analysed, this filter produced
non-degenerate patches for all 52.
\item \textbf{Sign-of-curvature calibration.} The sign of $H$ depends
on which direction is treated as ``outward.'' We anchor the convention
to a synthetic convex sphere through the same pipeline (see
\texttt{sphere\_sign\_check}); convex surfaces return $H > 0$. Any
break of this convention (e.g.~a degenerate cell whose mesh has
inverted normals) is detectable from the sphere check and was not
observed in production crops.
\item \textbf{Disconnected gap-channel components.} On crops where a
significant fraction of vertices are boundary-uncertain, the dropped
rim can physically separate the gap-channel submesh into multiple
connected components. This is rendered honestly --- the submesh
silhouette is irregular --- but should be interpreted with the
\texttt{bd-clip} fraction.
\end{itemize}

\subsection*{Visualization}
For the public gallery (\texttt{figures/membranes/membranes\_3d.html}),
the per-vertex scalars are vertex-coloured into glTF~2.0 (\texttt{.glb})
meshes via the \texttt{matplotlib} \texttt{RdBu\_r} (curvature,
protrusion) and \texttt{viridis} (gap) colormaps and rendered with
the \texttt{model-viewer} web component. Each card displays the
patch's three channels selectable in-place and also via a global
``color all by'' selector. The curvature and protrusion colormap
limits are symmetric and set to $\pm$\,p90 of $|H|$ and $|d|$ on the
patch (the brightest reds/blues are the sharpest features
\emph{in that patch}, not a fixed absolute threshold across crops).
The gap colormap is anchored at 0--$3 \cdot 40 = 120$\,nm (so crops
dominated by close contacts share a common scale) and expanded to
p95 of the kept-vertex gap distribution when the patch genuinely
sits hundreds of nm from any in-volume neighbour, with the active
upper bound shown on each card. Each card also exposes a
Neuroglancer link wired to the underlying crop with the EM, ECS
silhouette mesh, and all cell meshes pre-loaded; this allows
visual verification of any reading against the raw FIB-SEM EM.

\subsection*{Results --- overall}
% TODO: populate from manifest_3d.json once aggregate per-tissue
% summaries are computed. Suggested form: per-tissue table of
% (median $|H|$, IQR), (median $|d|$, IQR), (median $g$, IQR), and
% the bd-clip distribution, separated by prep. The matched-region
% comparison in the following subsection is the primary analytical
% result; this paragraph is a high-level descriptive aggregate.

\subsection*{Results --- region-matched comparisons}
% TODO: populate from manifest_3d.json + crop_annotations.csv. The
% matched region groups (e.g.~Liver Bile canaliculus, Liver Space of
% Disse, Liver Hepatocyte lateral; Heart Intercalated disc, Cardiac
% interstitial; Kidney PCT brush border, PCT lateral, DCT base,
% Glomerular interstitium; Cortex Neurovascular, Synaptic,
% Mitochondrion-rich) are the unit of comparison. Suggested form for
% each region group: paired per-channel (Chem vs HPF) per-vertex
% distributions, summarized as median + IQR, with a small inset
% mesh showing the most-representative crop.
